% !PRINTED none
% Abschnitt-opening page: no running-head folio number is printed on this
% leaf (suppressed on the chapter-opening page, as is common in this book);
% resumes as "564" on the next leaf.

\begin{center}
\gesperrt{Dreiundzwanzigster Abschnitt.}

{\Large \textbf{Der Teilungsk\"orper.}}
\end{center}

\begin{center}
\rule{2cm}{0.4pt}
\end{center}

\begin{center}
\S\,150. \textbf{Die homogenen \gesperrt{Weierstrass}schen Funktionen.}
\end{center}

Wir waren in \S\,114 von der Frage ausgegangen, unter welchen
Voraussetzungen eine doppelt periodische Funktion $\varphi(u)$ mit den
Perioden $\omega_1, \omega_2$ eine Multiplikation $\varphi(\mu u)$
gestattet, d.\,h. unter welchen Umst\"anden $\mu\omega_1, \mu\omega_2$
sich linear und ganzzahlig durch $\omega_1, \omega_2$ ausdr\"ucken
lassen, und waren zu dem Resultat gelangt, da\ss{} dies bei
ver\"anderlichem $\omega_1, \omega_2$ nur m\"oglich ist, wenn $\mu$ eine
ganze rationale Zahl ist, und da\ss{} nur, wenn

\weq{1}{\omega = \frac{\omega_2}{\omega_1}}

eine imagin\"are quadratische Irrationalzahl ist, auch $\mu$ eine
komplexe Zahl desselben K\"orpers wie $\omega$ sein kann. Es kommt
jetzt darauf an, die aus dieser Annahme folgenden Formeln der
komplexen Multiplikation etwas genauer zu erforschen.

Wir wollen zun\"achst, als das formal einfachere, die komplexe
Multiplikation der \gesperrt{Weierstrass}schen elliptischen Funktion $\wp(u)$
untersuchen, m\"ussen dann aber auch noch die komplexe Multiplikation
der \gesperrt{Jacobi}schen elliptischen Funktionen in Betracht ziehen.

In den \S\,37, (8), \S\,46, (13) haben wir die Formeln f\"ur die
Homogeneit\"at der Funktionen $\sigma(u)$, $\wp(u)$ und der
Invarianten $g_2, g_3, 16\,G = g_2^3 - 27\,g_3^2$ abgeleitet, die wir
hier noch einmal zusammenstellen:

\weq{2}{\begin{aligned}
\sigma(\lambda u, \lambda\omega_1, \lambda\omega_2) &= \lambda\ \sigma(u, \omega_1, \omega_2), \\
\wp(\lambda u, \lambda\omega_1, \lambda\omega_2) &= \lambda^{-2}\, \wp(u, \omega_1, \omega_2), \\
g_2(\lambda\omega_1, \lambda\omega_2) &= \lambda^{-4}\, g_2(\omega_1, \omega_2), \\
g_3(\lambda\omega_1, \lambda\omega_2) &= \lambda^{-6}\, g_3(\omega_1, \omega_2), \\
G(\lambda\omega_1, \lambda\omega_2) &= \lambda^{-12}\, G(\omega_1, \omega_2),
\end{aligned}}
%CONT>>
